\documentclass[10pt,a4paper]{amsart}
\usepackage{amsmath,amssymb,amsthm,graphicx}
\usepackage{hyperref}
\usepackage{tikz}
\usetikzlibrary{arrows}

\numberwithin{equation}{section}

\newtheorem{Theorem}{Theorem}[section]
\newtheorem*{Theorem*}{Theorem}
\newtheorem{Corollary}[Theorem]{Corollary}
\newtheorem{Lemma}[Theorem]{Lemma}
\newtheorem{Proposition}[Theorem]{Proposition}
\newtheorem{Conjecture}[Theorem]{Conjecture}

\theoremstyle{definition}
\newtheorem{Definition}[Theorem]{Definition}
\newtheorem{Note}[Theorem]{Note}
\newtheorem{Example}[Theorem]{Example}
\newtheorem{Remark}[Theorem]{Remark}
\usepackage{cleveref}
\crefname{originalreference}{original reference}{original references}
\Crefname{originalreference}{Original reference}{Original references}
\makeatletter
\newlabel{sl3-qm2}{{3.14}{source}{}{sourcepointer-1}{}}
\newlabel{sl3-qm2@cref}{{[originalreference][0][]3.14}{[1][0][]source}{}{}{}}
\makeatother
\title{Quasi-modularity of the sl3 Lie-theoretic q-MZV family via constant terms}
\author{Antun Milas}
\date{doi:10.3842/SIGMA.2026.087; CC BY 4.0}
\begin{document}
\maketitle
\noindent\textit{Allowed standard prerequisites (editorial).}\begin{itemize}
\item Standard q-binomial summation, vertex-algebra character definitions and classical elliptic/theta functions. Original source bibliography, including Lang, Arakawa and AM where relevant.
\end{itemize}\medskip
\ifdefined\hypertarget\hypertarget{sourcepointer-1}{}\fi
\noindent\textit{Source reference.} Original reference 3.14: algebra-geometry/sources/sigma-2026-087/native/sigma26-087.tex, source lines 1402--1402. Cross-reference gives an introductory specialization, earlier narrative comparison or section pointer; all source-local arguments used for the counted theorem are included.\par\medskip
Denote by $\Delta$ a root system of ADE type, by $\Delta_+$ the set of positive roots, and let $\langle \cdot, \cdot \rangle$ denote the inner product normalized
such that $\langle \alpha,\alpha \rangle=2$ for every root $\alpha$. For integers $k_\alpha \geq 1$, we define a $q$-MZV associated with $\mathfrak{g}$ as
\begin{equation} \label{qmzvg}
\zeta_{\mathfrak{g},q}(k_\alpha; \alpha \in \Delta_+):=\sum_{\lambda \in P_+} \frac{q^{\frac12 \sum_{\alpha \in \Delta_+} k_\alpha \langle\lambda+\rho,\alpha \rangle}}{\prod_{\alpha \in \Delta_+} (1-q^{\langle \lambda, \alpha+\rho\rangle})^{k_\alpha}},
\end{equation}
where the summation is over the cone of positive dominant integral weights. The asymptotic behavior motivates our definition. Indeed, after setting $q={\rm e}^{2\pi {\rm i}\tau}$ and taking the asymptotic limit $\tau \to 0+$ with
$k_{\alpha}:=k \geq 2$, we recover
Witten's $\zeta$-function of $\mathfrak{g}$ (see \cite{Zagier})
\[
\sum_{\lambda \in P_+} \frac{1}{\dim (V_\lambda)^k},
\]
 as the leading coefficient of the asymptotic expansion. This fact follows from the Weyl dimension formula.

In the simplest case of $\mathfrak{g}=\mathfrak{sl}_2$, where $\Delta=\{ \pm \alpha \}$, these $q$-zeta values are given by
\[
\zeta_{\mathfrak{sl}_2,q}(k)=\sum_{n \geq 1} \frac{q^{\frac{k}{2}n}}{(1-q^n)^k}.
\]
This slightly non-standard $q$-zeta value has several nice properties. For instance, the summand is invariant under $q \to q^{-1}$ for $k$ even and skew-invariant for $k$ odd. For simplicity, we only consider $k$ even here.
Then it is known (or it can be easily verified) that for $k \geq 1$, $\zeta_{\mathfrak{sl}_2,q}(2k)$, up to an additive constant, is a linear combination of Eisenstein series, and in fact we have
\begin{equation*} %\label{ordinary}
\mathbb{C}[\zeta_{\mathfrak{sl}_2,q}(2k): k \geq 1]=\underbrace{\mathbb{C}[E_2(q),E_4(q),E_6(q)]}_{:=\mathcal{QM}},
\end{equation*}
where $E_{2i}$ are Eisenstein series defined with their $q$-expansion; see formula (\ref{eis}) in Section~\ref{section3}.
The space $\mathcal{QM}$ is called the ring of (holomorphic) quasi-modular forms on ${\rm SL}(2,\mathbb{Z})$.

\section[Multiple q-zeta values, quasi-modular forms, and characters]{Multiple $\boldsymbol{q}$-zeta values, quasi-modular forms, and characters}\label{section3}

This section contains the second part of the paper. Throughout, we let
$\tau\in\mathbb H$, $z\in\mathbb C$, $q={\rm e}^{2\pi {\rm i}\tau}$, and
$\zeta={\rm e}^{2\pi {\rm i} z}$.
Unless stated otherwise, all identities involving Fourier series in $\zeta$
are interpreted in the annulus specified in the statement, which determines
the corresponding geometric-series expansion of each denominator. We begin by
recalling the elliptic and quasi-Jacobi functions needed for our constant-term
computations and then apply these techniques to characters of genus-zero
class-$\mathcal{S}$ vertex operator algebras.


\subsection{Elliptic functions, theta functions, and quasi-modular forms}\label{section3.1}
We present some basic (and mostly well-known) results on elliptic functions, (quasi)-Jacobi forms, and their Fourier coefficients. We closely follow \cite{Lang}.

The Jacobi theta function is given by
\[
\vartheta(z;\tau)=\sum_{n \in \mathbb{Z}+\frac12} {\rm e}^{\pi {\rm i} n^2 \tau+2 \pi {\rm i} n(z+\frac12)}.
\]
We also define Eisenstein series with their $q$-expansion, for integers $k\geq1$ (and with the usual convention that $E_2$ is quasimodular)
\[
G_{2k}(\tau)=2 \zeta(2k)+\frac{(2\pi {\rm i})^{2k}}{(2k-1)!} \sum_{n \geq 1} \frac{n^{2k-1}q^n}{1-q^n}, \qquad q={\rm e}^{2 \pi {\rm i} \tau},
\]
and slightly normalized versions
\begin{equation} \label{eis}
E_{2k}(\tau)=1-\frac{4k}{B_{2k}}\sum_{n \geq 1} \frac{n^{2k-1} q^n}{1-q^n},
\end{equation}
where $B_{2k}$ is the Bernoulli number. Recall that $\mathcal{QM}$ us the ring of quasi-modular forms for the full modular group discussed earlier. A famous result of S. Ramanujan is that
\[
\left(q \frac{{\rm d}}{{\rm d}q} \right)\colon \ \mathcal{QM} \rightarrow \mathcal{QM},
\]
so $\mathcal{QM}$ is a differential ring. This well-known fact is employed throughout this section and in later parts.

The Weierstrass $\zeta$-function\footnote{Not to be confused with the Fourier coefficient $\zeta={\rm e}^{2 \pi {\rm i} z}$ used below.} and the $\wp$-function are defined as
\begin{align*}
&\zeta(z;\tau) =\frac1z+\sum_{(n,m) \neq (0,0)} \left(\frac{1}{z-m \tau-n}+\frac{1}{m\tau+n}+\frac{z}{(m\tau+n)^2} \right),\\
&\wp(z; \tau) =\frac{1}{z^2}+\sum_{(n,m) \neq (0,0)} \frac{1}{(z-m-n\tau)^2}-\frac{1}{(m+n \tau)^2}, \\
&\wp'(z;\tau)=-\frac{1}{2} \frac{{\rm d}}{{\rm d}z} \wp(z;\tau).
\end{align*}
Their Laurent expansions around $z=0$ are then given by
\begin{align*}
&\zeta(z;\tau) =\frac{1}{z}-\sum_{n \geq 1} G_{2n+2}(\tau)z^{2n+1}, \\
&\wp(z;\tau)=\frac{1}{z^2} + \sum_{n \geq 1} (2n+1)G_{2n+2}(\tau) z^{2n}, \\
&\wp'(z;\tau)=\frac{1}{z^3} - \sum_{n \geq 1} n(2n+1) G_{2n+2}(\tau) z^{2n-1}.
\end{align*}
For concrete computation, it is convenient to use the so-called $P$-functions or ``propagators'' (widely used for computations of genus one correlators in CFT and vertex operator algebra~\cite{Zhu}):
\[
P_k(\zeta;q)=\frac{(2 \pi {\rm i})^k}{(k-1)!} \sum_{n \geq 1} \left( \frac{n^{k-1} \zeta^{n}q^n}{1-q^n}+ (-1)^k \frac{n^{k-1} \zeta^{-n}}{1-q^n} \right).
\]
Then inside $|q|<\bigl|\zeta^{-1}\bigr|<1$ we have
\begin{align}
\label{rho1} &P_1\bigl({\rm e}^{2 \pi {\rm i} z},\tau\bigr)=-\zeta(z;\tau)+G_2(\tau)z -\pi {\rm i}, \\
\label{rho2} &P_2\bigl({\rm e}^{2 \pi {\rm i} z}, \tau\bigr)=\wp(z,\tau)+G_2(\tau), \\
\label{rhok} &P_{3}\bigl({\rm e}^{2 \pi {\rm i} z}, \tau\bigr)=-\wp'(z,\tau).
\end{align}
Similarly, $P_k\bigl({\rm e}^{2 \pi {\rm i} z};\tau\bigr)$, $k \geq 4$, is proportional to the $(k-2)$-th $z$-derivative of $\wp(z;\tau)$.
The second and the third relations (\ref{rho2})--(\ref{rhok}) are equivalent to the following:
\begin{align*} %\label{rho-fourier}
& \wp(z,\tau)=-G_2(\tau)+(2 \pi {\rm i})^2 \sum_{n \geq 1} \left( \frac{n q^n \zeta^n }{(1-q^n)}+\frac{n \zeta^{-n} }{(1-q^n)} \right),\\
& \wp'(z,\tau)=-(2 \pi {\rm i})^3 \sum_{n \geq 1} \left( \frac{n^2 q^n \zeta^n }{(1-q^n)}-\frac{n^2 \zeta^{-n} }{(1-q^n)} \right).
\end{align*}
To see that (\ref{rho1}) holds, we require an identity
\begin{equation} \label{J1}
-\frac12 +\frac{\zeta}{\zeta-1} - \sum_{n \geq 1}\left(\frac{ \zeta^{-1} q^n}{1- \zeta^{-1}q^n}-\frac{q^n \zeta}{1-q^n \zeta}\right)=\frac12+\sum_{n \geq 1} \left(\frac{\zeta^n q^n}{1-q^n}-\frac{\zeta^{-n}}{1-q^n} \right),
\end{equation}
again valid in $|q|<|\zeta^{-1}|<1$. We denote either side of (\ref{J1}) by $J_1(z;\tau)$.
This function does not transform as a Jacobi form, instead we have
\begin{align*}% \label{J1-mod}
&J_1(z+ m \tau+n;\tau)= J_1(z;\tau)-m, \\
&J_1 \left(\frac{z}{\tau};\frac{-1}{\tau}\right)=\tau J_1(z;\tau)+z.
\end{align*}
We also require $\sigma(z;\tau)$, which can be defined using the theta function \cite{Lang}
\[
\sigma(z;\tau)=-{\rm exp}\left(\frac{\eta_1 z^2}{2}\right)\frac{\vartheta(z;\tau)}{2 \pi \eta(\tau)^3},
\]
where $\eta_1=G_2(\tau)$ in the normalization of \cite{Lang}; for our purposes, it is a function of $\tau$ only and will not be used elsewhere.

Functions $\wp$, $\vartheta$, and $\sigma$ are related through a fundamental relation~\cite{Lang}:
\begin{align}
\wp(z_1;\tau)-\wp(z_2;\tau)=\frac{\sigma(z_1+z_2;\tau)\sigma(z_1-z_2;\tau)}{\sigma(z_1;\tau)^2\sigma(z_2;\tau)^2} =\frac{\vartheta(z_1+z_2;\tau) \vartheta(z_1-z_2;\tau)}{ 4 \pi^2 \vartheta(z_1;\tau)^2 \vartheta(z_2;\tau)^2 \eta(\tau)^6}.\label{rho-sigma}
\end{align}

Our next result concerns the constant term in the Fourier expansion of elliptic functions with respect to $\zeta = {\rm e}^{2\pi {\rm i} z}$.
This result is well known and has appeared previously in the literature; see, for example, \cite{GM, MP}.
In the present work, however, we adopt the approach of~\cite{BFM}, where all Fourier coefficients of certain index-zero Jacobi forms are investigated, rather than focusing solely on the constant term.
\begin{Proposition} \label{elliptic} Let $f(z;\tau) \in \mathbb{C}[\wp(z;\tau),\wp'(z;\tau)]$.
Then, in the range $ |q|< \bigl|\zeta^{-1}\bigr|<1$, the constant term ${\rm CT}_\zeta f(z;\tau)$ belongs to $\mathcal{QM}$, and the non-constant terms ${\rm CT}_\zeta f(z;\tau)\zeta^{-r}$, $r \neq 0$, are not quasi-modular.
\end{Proposition}
\begin{proof} We may assume that $f=\wp^r (\wp')^s$. Using the Weierstrass equation
\[
(\wp')^2=4\wp^3-g_2(\tau)\wp-g_3(\tau),
\]
where $g_2(\tau)$ and $g_3(\tau)$ are multiples of $E_4(\tau)$ and $E_6(\tau)$, respectively, every higher power of $\wp'$ can be reduced to the cases $s=0$ and $s=1$.
We use the Laurent expansion in the variable $z$ around zero to write
\begin{align*} %\label{formula}
f(z;\tau)=\sum_{0 \leq i < \frac{\ell}{2}} \frac{D_{\ell-2i}(\tau)}{ z^{\ell-2i}}+D_0(z)+O(z),
\end{align*}
where the $z$-coefficients $D_{\ell-2i}(\tau)$ are holomorphic {modular} forms of weight $2i$ that can be explicitly computed using the Laurent expansion of $\wp$
and $\wp'$, and $\ell=2r+3s$. Of course, $D_0(z)$ could be zero depending on the parity of $s$. To compute the $\zeta$-constant term, we proceed as in \cite[Theorem~1.1]{BFM} dealing with a particular Jacobi form of index zero (the case when $m$ is even in~\cite{BFM} is precisely the case of elliptic functions here).

For the nonzero Fourier coefficient $f_r(\tau)$, $r \neq 0$, we compute the contour integral
\[
\int_{\partial P_{z_0}} f(w;\tau) {\rm e}^{-2 \pi {\rm i} r w} \, {\rm d}w,
\]
where
\[
P_{z_0} = \{ z_0 + r \tau + s \mid 0 \leq r,s \leq 1 \},
\]
in two different ways: directly as a contour integral, and by the residue theorem by choosing
\[
z_0 = -\tfrac{1}{2} - \tfrac{\tau}{2},
\]
so that the only pole inside the contour is at $z=0$.
This quickly leads to a formula for the $r$-th coefficient, which is quasi-modular if and only if $r=0$ and not quasi-modular otherwise (see \cite[Theorem~1.1]{BFM}).
More explicitly,
\[
{\rm CT}_\zeta f(z;\tau)=D_0(\tau)+\sum_{i} \lambda_i D_{2i}(\tau)E_{2i}(\tau),
\]
where $\lambda_i$ are some coefficients. This is certainly contained in~$\mathcal{QM}$.
\end{proof}

\subsection{Quasi-Jacobi forms}\label{section3.2}
We are also interested in Fourier coefficients, or more precisely the constant term, of more general functions that include $\tau$-derivatives of elliptic
functions. For that we need the concept of {\em quasi-Jacobi forms} (see \cite{Libgober,Ob, Ittersum} for instance).
We are only concerned with the subspace of meromorphic quasi-Jacobi forms of index $0$ with poles at $L_\tau=\mathbb{Z}\tau+\mathbb{Z}$. These are known to form a subring $\mathcal{QJ}^0:=\mathbb{C}[E_6,E_4,J_1,\wp,\wp']$ inside the ring of quasi-Jacobi forms $\mathcal{QJ}:=\mathbb{C}[F,E_6,E_4,J_1,\wp,\wp']$ ($E_6$ can be omitted from both rings though), where $F$ is a particular quasi-Jacobi form of non-zero index not needed here \cite{Ob}. Besides $J_1(z;\tau)$ mentioned earlier, we also require more general quasi-Jacobi forms ($k \geq 2$):
\[
J_k(z;\tau)=B_k - k \sum_{n, r \geq 1} r^{k-1}\bigl( \zeta^n + (-1)^k \zeta^{-n}\bigr) q^{nr}.
\]
Next we define the polynomials $p_m(x)$, $m \geq 2$ via
\[
\sum_{n \geq 1} n^{m-1} x^n=\frac{x p_{m}(x)}{(1-x)^{m}}.
\]
For instance, $p_2(x)=1$, $p_3(x)=1+x$, etc.
These are known as (slightly shifted) Euler's polynomials. The next result is elementary and essentially follows from the definition of $p_m(x)$.
\begin{Lemma} \label{euler} We have
\begin{itemize}\itemsep=0pt
\item[$(1)$] For every even $s \geq 0$, and $k \geq 1$,
\[
\sum_{r \geq 1} \frac{r^s q^r p_{2k}(q^r)}{(1-q^r)^{2k}} \in \mathcal{QM}.
\]
\item[$(2)$] For every odd $s \geq 1$, and $k \geq 1$,
\[
\sum_{r \geq 1} \frac{r^s q^r (1+q^r) p_{2k}(q^r)}{(1-q^r)^{2k+1}} \in \mathcal{QM}.
\]
\end{itemize}
\end{Lemma}
\begin{proof} For (1), it is enough to observe:
\begin{align*}
 \sum_{r \geq 1} r^j \frac{q^r p_{a+j}(q^r)}{(1-q^r)^{a+j}}= \sum_{r \geq 1,\, n \geq 1}\! r^j n^{a+j-1} q^{n r}= \left(q \frac{{\rm d}}{{\rm d}q} \right)^j \! \sum_{r \geq 1,\, n \geq 1}\! n^{a-1} q^{n r} =\left(q \frac{{\rm d}}{{\rm d}q}\right)^j\! \sum_{n \geq 1} \frac{n^{a-1} q^n}{1-q^n}.
\end{align*}
This implies that, for every $s$ and $k$ as above, we may choose $a \geq 2$ and $j \geq 0$ such that the sum in (1) is a (higher) derivative of an Eisenstein series, or an Eisenstein series with its constant term removed.
Similarly, any series appearing in (2) can be expressed as a linear combination of derivatives of two Eisenstein series.
In either case, the resulting sum lies in $\mathcal{QM}$.
\end{proof}

\begin{Theorem} \label{quasi} Let $f(z;\tau) \in \mathcal{QJ}^0$. Then in the range $|q|<\bigl|\zeta^{-1}\bigr|<1$, ${\rm CT}_\zeta f(z;\tau) \in \mathcal{QM}$.
\end{Theorem}
\begin{proof}
Compared with Proposition~\ref{elliptic}, where $f$ is doubly periodic in $z$ and belongs to $\mathbb{C}[\wp,\wp']$, the only new issue is the contribution of $J_1(z;\tau)$ and thus it suffices to prove the statement for
\begin{align*}
f(z;\tau)=J_1(z;\tau)^j g(z;\tau),
\end{align*}
where $g(z;\tau) \in \mathbb{C}[\wp,\wp']$ is elliptic and $j \geq 1$.
We use induction on $j$. For $j=1$, we have to
analyze
\[
J_1(z;\tau) g(z;\tau),
\]
where $g(z;\tau)$ is elliptic. We will compute the constant term by expanding $g(z;\tau)=\zeta^r g_r(\tau)$, where $g_0(\tau)$ is quasi-modular (possibly zero) and
non-zero Fourier coefficients are given by~\cite{BFM}
\[
g_r(\tau)=\frac{q^{r}}{1-q^r} \sum_{0 \leq \ell < m/2} r^{m-2\ell-1} D_{\ell}(\tau), \qquad r>0,
\]
and
\[
g_r(\tau)=-\frac{1}{1-q^{-r}} \sum_{0 \leq \ell < m/2} r^{m-2\ell-1} D_{\ell}(\tau), \qquad r<0,
\]
where the parity of $m$ depends on the parity of the $z$ powers in the Laurent expansion of $g(z;\tau)$, and $D_{\ell}(\tau)$ are modular forms as in Proposition~\ref{elliptic}.
Multiplying out $J_1(z;\tau)=\frac12+\sum_{n \geq 1} \bigl(\frac{\zeta^n q^n}{1-q^n}-\frac{\zeta^{-n}}{1-q^n} \bigr)$ and $g(z;\tau)$, we quickly find the constant term of $J_1(z;\tau) g(z;\tau)$ to be
\[
\frac{1}{2}g_0(\tau)- \sum_{0 \leq \ell < m/2} D_\ell(\tau) \sum_{r \geq 1} \frac{ r^{m-2\ell-1} q^r}{(1-q^r)^2} - \sum_{0 \leq \ell < m/2} D_\ell(\tau) \sum_{r \geq 1} \frac{ (-r)^{m-2\ell-1} q^r}{(1-q^r)^2}.
\]
If $m$ is even, then the two sums cancel out and we end up with $\frac12 g_0(\tau)$ for the constant term which is certainly in $\mathcal{QM}$ by Proposition~\ref{elliptic}. If $m$ is odd, then we get a contribution from the sum
\[
\sum_{r \geq 1} \frac{ r^{s} q^r}{(1-q^r)^2},
\]
where $s:=m- 2\ell -1 \geq 0$ is even. For $s=0$, this sum equals $\frac{1}{24}(1-E_2(\tau)) \in \mathcal{QM}$. For $s \geq 2$, the sum is the first $\bigl(q \frac{{\rm d}}{{\rm d}q}\bigr)$-derivative of
the Eisenstein series $E_{s-1}(\tau)$ and thus in $\mathcal{QM}$.

For the induction step, we need \cite[formula~(11)]{Ob-Serre} that allows us to express $J_1^\ell$ as a linear combination of $J_\ell$, an elliptic Jacobi form $K_\ell$, and quasi-Jacobi forms involving $J_1^k$, $k < \ell$, for which the induction hypothesis applies.
Thus in order to finish the proof we have to argue that
\[
{\rm CT}_\zeta J_\ell(z;\tau) g(z;\tau) \in \mathcal{QM},
\]
where $g(z;\tau)$ is elliptic. As in the $j=1$ case, we first observe that
\[
\frac{1}{\ell}(B_{\ell}-J_{\ell}(z;\tau))=\sum_{n \geq 1}\left( \frac{q^n p_\ell(q^n) \zeta^n}{(1-q^n)^\ell} + (-1)^\ell \frac{q^n p_\ell(q^n) \zeta^{-n}}{(1-q^n)^\ell} \right),
\]
where $p_\ell$ are as above. We distinguish two cases: $\ell$ is even and $\ell$ is odd.
Assume first that $\ell=2k$ is even.
After we multiply out Fourier expansions and extract the constant term as above, the constant term takes the shape
\[
\sum_{r \geq 1} \frac{q^{2r} r^s p_{2k}(q^r)}{(1-q^r)^{2k+1}} - (-1)^{s} \sum_{r \geq 1} \frac{q^{r} r^s p_{2k}(q^r)}{(1-q^r)^{2k+1}},
\]
where $s:=m - 2 n-1$ could be either even or odd. For $s$ even, we have
\begin{equation*} %\label{even1}
\sum_{r \geq 1} \frac{r^s q^{2r} p_{2k}(q^r)}{(1-q^r)^{2k+1}} - (-1)^{s} \sum_{r \geq 1} \frac{r^s q^{r} p_{2k}(q^r)}{(1-q^r)^{2k+1}}=- \sum_{r \geq 1} \frac{r^s q^{r}p_{2k}(q^r)}{(1-q^r)^{2k}}
\end{equation*}
whereas for $s$ odd, we obtain
\begin{equation*} %\label{even2}
\sum_{r \geq 1} \frac{r^s q^{2r} p_{2k} (q^r)}{(1-q^r)^{2k+1}} + \sum_{r \geq 1} \frac{r^s q^{r} p_{2k}(q^r)}{(1-q^r)^{2k+1}}=\sum_{r \geq 1} \frac{r^s q^{r}(1+q^r) p_{2k}(q^r)}{(1-q^r)^{2k+1}}.
\end{equation*}
By Lemma \ref{euler}, both series are in $\mathcal{QM}$.
Similarly we argue for $\ell$ odd. That completes the proof.
\end{proof}

Denote by $D_z=\frac{1}{2 \pi {\rm i} } \frac{\partial}{\partial z}$ and ${D}_\tau=\frac{1}{2 \pi {\rm i}} \frac{\partial}{\partial \tau},$ and let $s(x,y) \in \mathbb{C}[x,y]$ be a polynomial with complex coefficients.
\begin{Corollary} \label{const} Let $f(z;\tau) \in \mathcal{QJ}^{0}$. Then
\[
{\rm CT}_\zeta (s(D_z,D_\tau) f(z;\tau) ) \in \mathcal{QM}.
\]
\end{Corollary}
\begin{proof} We only have to argue that the ring $\mathcal{QJ}^{0}$ is closed under $D_z$ and $D_\tau$. This can be verified explicitly on the generators $J_1$, $\wp$, $\wp'$, $E_4$, $E_6$ (see, for instance, \cite{Ob, Ittersum}).
\end{proof}

An important special case involves products of derivatives of the $P$-functions introduced earlier.
A similar result appeared previously in \cite[Proposition~5.9]{GM}, though our proof differs.

\begin{Proposition} \label{P-power} We have
\begin{itemize}\itemsep=0pt
\item[$(i)$] For $m \geq 1$ and $k \geq 1$,
\[
{\rm CT}_{\zeta} \bigl(P_k(\zeta;\tau)^m \bigr) \in \mathcal{QM}.
\]
\item[$(ii)$] More generally, for every polynomial $s(x,y)$ and $k \geq 1$,
\[
{\rm CT}_{\zeta} \bigl\{ (s(D_\tau,D_{z}) P_k(\zeta;\tau))^{m} \bigr\} \in \mathcal{QM}.
\]
\end{itemize}
\end{Proposition}
\begin{proof} Part (ii) follows immediately from Theorem~\ref{quasi}, Corollary \ref{const}, and formulas (\ref{rho1})--(\ref{rhok}).
\end{proof}

\subsection[Characters of genus zero class-S vertex operator algebras]{Characters of genus zero class-$\boldsymbol{\mathcal{S}}$ vertex operator algebras}\label{section3.3}

In this section, $\mathfrak{g}$ denotes a finite dimensional Lie algebra of ADE type for simplicity. In the introduction, we
constructed a $q$-MZV function associated with $\mathfrak{g}$, $\zeta_{\mathfrak{g},q}(k_{\alpha}; \alpha \in \Delta_+)$. Letting $k_{\alpha} \geq 1$, we consider a slightly different $q$-MZV model where all $k_{\alpha}$ are equal and
there is an extra polynomial dependence on the labels of $\lambda \in P_+$:
\[
\zeta_{\mathfrak{g},q}^{s}(k):=\sum_{\lambda \in P_+} (\dim (V(\lambda))^{s} \frac{q^{k \langle\lambda+\rho,\rho\rangle}}{\prod_{\alpha \in \Delta_+} (1-q^{\langle \lambda, \alpha+\rho\rangle})^{k}}.
\]
The boldface symbol $\mathbf{V}_{\mathfrak g,k}$ denotes Arakawa's vertex operator algebra attached to the Lie algebra~$\mathfrak g$ (actually our notation $\mathbf V_{\mathfrak g,k}$ corresponds to Arakawa's $\mathbf V^{\mathcal S}_{G,b}$ with $\mathfrak g=
\operatorname{Lie}(G)$ and $b=k+2$; see~\cite{Arakawa})
and the positive integer parameter~$k$; it should not be confused with the ordinary italic letter $k$ used as a summation or weight parameter elsewhere.
Here $V(\lambda)$ is the finite-dimensional $\mathfrak{g}$-module of highest weight $\lambda$, so that $\dim (V(\lambda))$ can be computed using the Weyl dimension formula and $\rho=\frac12 \sum_{\alpha \in \Delta_+} \alpha$ is the half-sum of positive roots. Then we have a remarkable result of Arakawa on class-$\mathcal{S}$ vertex operator algebras motivated by the previous very important work of Beem, Rastelli and others~\cite{BR,Beem}.

\begin{Theorem}[\cite{Arakawa}] For every $k \geq 3$, there is a vertex operator algebra ${\bf V}_{\mathfrak{g},k}$ such that
\begin{equation} \label{Ar-char}
{\rm ch}[{\bf V}_{\mathfrak{g},k}](\tau)=\frac{1}{\eta(\tau)^{m_{\mathfrak{g},k}}} \zeta_{\mathfrak{g},q}^{k}(k-2),
\end{equation}
where $m_{\mathfrak{g},k}=\dim (\mathfrak{g})k-(k-2){\rm rank}(\mathfrak{g})$ is a positive integer and $\eta(\tau)=q^{1/24}\prod_{i \geq 1}(1-q^i)$ is the Dedekind $\eta$ function.
\end{Theorem}

\begin{Remark}
Although this seems to be implicit in the literature, it is expected that
\[
\dim X_{\mathbf V_{\mathfrak g,k}}=m_{\mathfrak g,k},
\]
where $X_{\mathbf V_{\mathfrak g,k}}$ is the associated variety of $\mathbf V_{\mathfrak g,k}$. In the 4d/2d correspondence, the associated variety is expected to recover the Higgs branch of the corresponding four-dimensional theory, while the vertex operator algebra character is identified with its Schur index. Thus formula~(\ref{Ar-char}) suggests that the generalized $q$-MZV is the infinite-dimensional analogue of the numerator in a Hilbert-series presentation: it contains the non-free contribution after the universal product factor has been separated. We use the phrase ``Hilbert-series numerator'' here, rather than ``$h$-series'', to avoid confusion with the $h$-polynomial of a finitely generated graded algebra.
\end{Remark}

\begin{Example}[$\mathfrak{g}=\mathfrak{sl}_2$] For $k \geq 3$,
\[
\zeta_{\mathfrak{sl}_2,q}^k(k-2)=\sum_{n \geq 1} \frac{n^k q^{\frac{n}{2} (k-2)} }{(1-q^n)^{k-2}}.
\]
\end{Example}

\begin{Example}[$\mathfrak{g}=\mathfrak{sl}_3$] For $k \geq 3$,
\[
\zeta_{\mathfrak{sl}_3,q}^k(k-2)=\sum_{m_1,m_2 \geq 1} \left( m_1 m_2 (m_1+m_2) \right)^k \frac{q^{(m_1 +m_2)(k-2)}}{(1-q^{m_1})^{k-2} (1-q^{m_2})^{k-2}(1-q^{m_1+m_2})^{k-2}}.
\]
\end{Example}

It is not hard to see that we also have
\begin{align*} %\label{type-An}
 \zeta_{\mathfrak{sl}_{n+1},q}^k(k-2)
 =\sum_{m_1,\dots,m_n \geq 1} \left( \prod_{1 \leq i \leq j \leq n} (m_i+\cdots + m_j)^k\right) \cdot \frac{q^{(k-2)( \frac12 \sum_{i=1}^n (i-1+1)(n-i+1)m_i)}}{\prod_{1 \leq i \leq j \leq n} (1-q^{m_i+\cdots + m_j})^{k-2}}.
\end{align*}
This can be rewritten as
\begin{align*} %\label{type-An-alt}
 \zeta_{\mathfrak{sl}_{n+1},q}^k(k-2) =\sum_{m_1,\dots,m_n \geq 1} \prod_{1 \leq i \leq j \leq n} \left( (m_i+\cdots + m_j)^k \frac{q^{\frac12(k-2)( i m_i+\cdots + j m_j)}}{(1-q^{m_i+\cdots + m_j})^{k-2}}\right).
\end{align*}

Based on the $\mathfrak{sl}_2$ case studied by Beem and Rastelli \cite{Beem}, it is expected the following.

\begin{Conjecture} \label{MLDE} We have
\begin{itemize}\itemsep=0pt
\item[$(1)$] For $k \geq 3$ even, $\zeta_{\mathfrak{sl}_{n+1},q}^k(k-2)$ belongs to $\mathcal{QM}$ and for $k \geq 3$ odd it belongs to $\mathcal{QM}[2]:=\mathbb{Q}\bigl[E_2(\tau), \theta_2(\tau)^4,\theta_3(\tau)^4\bigr]$.\footnote{It is possible to formulate a slightly stronger conjecture depending on whether or not $\rho \in Q$, where $Q$ is the root lattice of~$\mathfrak{g}$.}
\item[$(2)$]For $k \geq 3$, the character ${\rm ch}[\mathbf{V}_{\mathfrak{g},k}](\tau)$ is a solution of a modular linear differential equation $($MLDE$)$ with coefficients in $\mathcal{QM}$ or $\mathcal{QM}[2]$, depending on the parity of~$k$.
\end{itemize}
\end{Conjecture}
Beem and Rastelli verified these conjectures for numerous examples for $\mathfrak{g}=\mathfrak{sl}_2$ (for both $k$ even and odd) and also for additional examples coming from
genus $g \geq 1$ theories. These higher-genus theories are not discussed in Arakawa's paper but presumably correspond to characters of (mostly unknown) vertex operator algebras due to 2d/4d correspondences \cite{BR}. In \cite{BVM}, we proved a broader theorem that implies the even-parameter quasi-modularity assertion in Conjecture~\ref{MLDE}\,(1) for the relevant $A$-type series. It does not by itself prove the MLDE assertion in Conjecture~\ref{MLDE}\,(2). Theorem~\ref{sl3-qm} below supplies a direct, self-contained constant-term proof for the first non-trivial case $\mathfrak{g}=\mathfrak{sl}_3$, while Proposition~\ref{sl3-qm2} treats the corresponding symmetrized family.

\subsection[Bi-brackets of multiple q-zeta values of a Lie algebra g]{Bi-brackets of multiple $\boldsymbol{q}$-zeta values of a Lie algebra $\boldsymbol{\mathfrak{g}}$}\label{section3.4}

The notation in this subsection is independent of the one-variable zeta values in (\ref{qmzvg}); here the square-bracket notation records two arrays of parameters, following the terminology of bi-brackets.
Motivated by Arakawa's character formula, we can also define more general multiple $q$-zeta series, reminiscent of the {\em bi-brackets} construction in the setup of $q$-MZVs \cite{BK2,BK1,BK3}. For simplicity, we again only discuss the $A$-type here so the positive roots are parametrized by the segments $[i,j]$, $i \leq j$.

Let $k_{[i,j]} \geq 1$, $1 \leq i \leq j$ and $s_{[i,j]} \geq 0$ be integers, and
\begin{align*}
 & \zeta_{\mathfrak{sl}_{n+1},q} \left[ \begin{matrix} k_{[i,j]}; & 1 \leq i \leq j \leq n \\ s_{[i,j]}; & 1 \leq i \leq j \leq n \end{matrix} \right]\nonumber\\
 &\qquad :=\sum_{m_1,\dots,m_n \geq 1} \ \prod_{1 \leq i \leq j \leq n} (m_i+\cdots + m_j)^{s_{[i,j]}} \frac{q^{\frac12 k_{[i,j]} (i m_i+\cdots + j m_j)}}{(1-q^{m_i+\cdots + m_j})^{k_{[i,j]}}}. \label{type-An-alt}
\end{align*}
With this definition, we have
\[
\zeta_{\mathfrak{sl}_{n+1},q} \left[ \begin{matrix} k_{[i,j]} =k; & 1 \leq i \leq j \leq n \\ s_{[i,j]} =k+2; & 1 \leq i \leq j \leq n \end{matrix} \right]
=\zeta_{\mathfrak{sl}_{n+1},q}^{k+2} (k).
\]
Based on many examples, we expect a symmetric formula to hold.

\begin{Conjecture} Assume that all $k_{[i,j]}$ are even and $s=s_{[i,j]}$ for all $[i,j]$. Then
\[
\sum_{\sigma \in W} \zeta_{\mathfrak{sl}_{n+1},q} \left[ k_{\sigma[1,1]},k_{\sigma[1,2]}, \dots, k_{\sigma[n,n]};s \right] \in \mathcal{QM},
\]
where $k_{-[i,j]}=k_{[i,j]}$.
\end{Conjecture}

One can also formulate a similar conjecture for other root systems.
We also expect quasi-modularity to hold for
$\zeta_{\mathfrak{sl}_{n+1},q} \left[ \begin{smallmatrix} k_{[i,j]}; & 1 \leq i \leq j \leq n \\ s_{[i,j]}; & 1 \leq i \leq j \leq n \end{smallmatrix} \right]$
provided that $k_{\alpha}$ and $s_{\alpha}$ have the same parity.


\subsection[Example: Quasi-modularity for sl(3)]{Example: Quasi-modularity for $\boldsymbol{\mathfrak{sl}(3)}$}\label{section3.5}

The quasi-modularity of $\mathfrak{sl}(2)$ $q$-MZV values and their modifications is well known to experts, and therefore we only briefly comment on this case.
Recall the Lie-theoretic $q$-multiple zeta values introduced in
\eqref{qmzvg}. In the simplest case $\mathfrak g=\mathfrak{sl}_2$, one has
\[
\zeta_{\mathfrak{sl}_2,q}(2k)
=
\sum_{n\ge1}\frac{q^{nk}}{(1-q^n)^{2k}}.
\]
It is well known that, up to an additive constant, these Lambert series are
linear combinations of Eisenstein series and consequently,
\[
\zeta_{\mathfrak{sl}_2,q}(2k)\in\mathcal{QM}.
\]
Similarly, for every $k\geq1$, there exist constants
$c_{k,i}\in\mathbb{Q}$, $1\leq i\leq k$, such that, for every
$n\geq1$ and $a\geq0$, the following identity of rational functions in $q$
holds:
\begin{equation}\label{lin-comb}
\frac{n^{2k+a}q^{nk}}{(1-q^n)^{2k}}
=
\sum_{i=1}^{k}
c_{k,i}
\left(q\frac{{\rm d}}{{\rm d}q}\right)^{2k+1-2i}
\frac{n^{2i-1+a}q^n}{1-q^n}.
\end{equation}

Summing over $n$ and letting $a=2$ immediately gives
\[
\zeta_{\mathfrak{sl}_2,q}^{2k+2}(2k)=\sum_{i=1}^{k} c_{k,i} \left(q \frac{{\rm d}}{{\rm d}q}\right)^{2k+1-2i} E_{2+2i}(q) \in \mathbb{Q}[E_2,E_4,E_6].
\]

\begin{Corollary}
For integers $k\geq1$ and $a\geq0$, we have
\begin{align}
&\sum_{i=1}^{k} c_{k,i}
D_\tau^{\,2k+1-2i}
\left(
\sum_{n\geq1}
\left(
\frac{n^{2i-1+a}q^n\zeta^n}{1-q^n}
+
(-1)^{a+1}
\frac{n^{2i-1+a}q^n\zeta^{-n}}{1-q^n}
\right)
\right)\nonumber
\\
&\qquad =
\sum_{n\geq1}
\left(
\frac{n^{2k+a}q^{nk}\zeta^n}{(1-q^n)^{2k}}
+
(-1)^a
\frac{n^{2k+a}q^{nk}\zeta^{-n}}{(1-q^n)^{2k}}
\right).\label{form-sl3}
\end{align}
\end{Corollary}
Equipped with (\ref{form-sl3}) and results from Section~\ref{section3.3}, we are now ready to prove the following.


\begin{Theorem} \label{sl3-qm} Let $k\geq1$ and let $a\geq0$ be even. Then
\[
\sum_{m_1,m_2\geq1}\left(m_1m_2(m_1+m_2)\right)^{2k+a}
\frac{q^{2k(m_1+m_2)}}{(1-q^{m_1})^{2k}(1-q^{m_2})^{2k}(1-q^{m_1+m_2})^{2k}}
\in\mathcal{QM}.
\]
In particular, taking $a=2$ gives $\zeta_{\mathfrak{sl}_3,q}^{2k+2}(2k)\in\mathcal{QM}$.
\end{Theorem}

\begin{proof} Keeping $a$ even, we prove that
\[
\zeta_{\mathfrak{sl}_3,q}^{2k+a}(2k)=\sum_{m_1,m_2 \geq 1} \left( m_1 m_2 (m_1+m_2) \right)^{2k+a} \frac{q^{2k(m_1 +m_2)}}{(1-q^{m_1})^{2k} (1-q^{m_2})^{2k}(1-q^{m_1+m_2})^{2k}}
\]
also belongs to $\mathcal{QM}$.
Keeping $a$ even, observe an identity
\begin{align*}
& {\rm CT}_{\zeta} \left(\sum_{n \geq 1} \frac{n^{2k+a} q^{nk} \zeta^{n}}{(1-q^n)^{2k}}+(-1)^a \frac{n^{2k+a} q^{n k} \zeta^{-n}}{(1-q^n)^{2k}} \right)^3 \\
&\qquad = {\rm CT}_{\zeta} \left(\sum_{n \geq 1} \frac{n^{2k+a} q^{nk} \zeta^{n}}{(1-q^n)^{2k}}+ \frac{n^{2k+a} q^{n k} \zeta^{-n}}{(1-q^n)^{2k}} \right)^3 \\
& \qquad =6 \sum_{m_1,m_2 \geq 1} \left( m_1 m_2 (m_1+m_2) \right)^{2k+a} \frac{q^{2k(m_1 +m_2)}}{(1-q^{m_1})^{2k} (1-q^{m_2})^{2k}(1-q^{m_1+m_2})^{2k}},
\end{align*}
where in the last step we expanded the triple product as triple sums and we extract the constant term from each individual term.
Formula~(\ref{form-sl3}), together with Proposition~\ref{P-power}, now gives the assertion.
\end{proof}

\begin{thebibliography}{99}
\footnotesize\itemsep=0pt

\bibitem{AM}
Adamovi\'{c} D., Milas A., On some vertex algebras related to
 {$V_{-1}(\mathfrak{sl}(n))$} and their characters,
 \href{https://doi.org/10.1007/s00031-020-09617-w}{\textit{Transform. Groups}}
 \textbf{26} (2021), 1--30,
 \href{https://arxiv.org/abs/1805.09771}{arXiv:1805.09771}.

\bibitem{Andrews0}
Andrews G.E., The theory of partitions, \textit{Encyclopedia Math. Appl.},
 Vol.~2, Addison-Wesley Publishing Co., Reading, Mass., 1976.

\bibitem{Andrews}
Andrews G.E., Stacked lattice boxes,
 \href{https://doi.org/10.1007/BF01608779}{\textit{Ann. Comb.}} \textbf{3}
 (1999), 115--130.

\bibitem{AR}
Andrews G.E., Rose S.C.F., Mac{M}ahon's sum-of-divisors functions, {C}hebyshev
 polynomials, and quasi-modular forms,
 \href{https://doi.org/10.1515/crelle.2011.179}{\textit{J.~Reine Angew.
 Math.}} \textbf{676} (2013), 97--103,
 \href{https://arxiv.org/abs/1010.5769}{arXiv:1010.5769}.

\bibitem{Arakawa}
Arakawa T., Chiral algebras of class $\mathcal{S}$ and {M}oore--{T}achikawa
 symplectic varieties,
 \href{https://arxiv.org/abs/1811.01577}{arXiv:1811.01577}.

\bibitem{AK}
Arakawa T., Kawasetsu K., Quasi-lisse vertex algebras and modular linear
 differential equations, in Lie Groups, Geometry, and Representation Theory,
 \textit{Progr. Math.}, Vol. 326,
 \href{https://doi.org/10.1007/978-3-030-02191-7_2}{Birkh\"auser}, Cham, 2018,
 41--57, \href{https://arxiv.org/abs/1610.05865}{arXiv:1610.05865}.

\bibitem{BK2}
Bachmann H., K\"uhn U., A short note on a conjecture of {O}kounkov about a
 $q$-analogue of multiple zeta values,
 \href{https://arxiv.org/abs/1407.6796}{arXiv:1407.6796}.

\bibitem{BK1}
Bachmann H., K\"uhn U., The algebra of generating functions for multiple
 divisor sums and applications to multiple zeta values,
 \href{https://doi.org/10.1007/s11139-015-9707-7}{\textit{Ramanujan~J.}}
 \textbf{40} (2016), 605--648,
 \href{https://arxiv.org/abs/1309.3920}{arXiv:1309.3920}.

\bibitem{BK3}
Bachmann H., K\"uhn U., A dimension conjecture for $q$-analogues of multiple
 zeta values, in Periods in Quantum Field Theory and Arithmetic,
 \textit{Springer Proc. Math. Stat.}, Vol.~314,
 \href{https://doi.org/10.1007/978-3-030-37031-2_9}{Springer}, Cham, 2020,
 237--258, \href{https://arxiv.org/abs/1708.07464}{arXiv:1708.07464}.

\bibitem{BR}
Beem C., Lemos M., Liendo P., Peelaers W., Rastelli L., van Rees B.C., Infinite
 chiral symmetry in four dimensions,
 \href{https://doi.org/10.1007/s00220-014-2272-x}{\textit{Comm. Math. Phys.}}
 \textbf{336} (2015), 1359--1433,
 \href{https://arxiv.org/abs/1312.5344}{arXiv:1312.5344}.

\bibitem{Beem}
Beem C., Rastelli L., Vertex operator algebras, {H}iggs branches, and modular
 differential equations,
 \href{https://doi.org/10.1007/jhep08(2018)114}{\textit{J.~High Energy Phys.}}
 \textbf{2018} (2018), no.~8, 114, 72~pages,
 \href{https://arxiv.org/abs/1707.07679}{arXiv:1707.07679}.

\bibitem{Beem-new}
Beem C., Singh P., Razamat S.S., Schur indices of class {$\mathcal{S}$} and
 quasimodular forms,
 \href{https://doi.org/10.1103/physrevd.105.085009}{\textit{Phys. Rev.~D}}
 \textbf{105} (2022), 085009, 15~pages,
 \href{https://arxiv.org/abs/2112.10715}{arXiv:2112.10715}.

\bibitem{Bradley}
Bradley D.M., Multiple {$q$}-zeta values,
 \href{https://doi.org/10.1016/j.jalgebra.2004.09.017}{\textit{J.~Algebra}}
 \textbf{283} (2005), 752--798,
 \href{https://arxiv.org/abs/math/0402093}{arXiv:math/0402093}.

\bibitem{BFM}
Bringmann K., Folsom A., Mahlburg K., Quasimodular forms and
 {$s\ell(m|m)^\wedge$} characters,
 \href{https://doi.org/10.1007/s11139-014-9621-4}{\textit{Ramanujan~J.}}
 \textbf{36} (2015), 103--116,
 \href{https://arxiv.org/abs/1302.4040}{arXiv:1302.4040}.

\bibitem{BJM}
Bringmann K., Jennings-Shaffer C., Milas A., Graph schemes, graph series, and
 modularity,
 \href{https://doi.org/10.1016/j.jcta.2023.105749}{\textit{J.~Combin. Theory
 Ser.~A}} \textbf{197} (2023), 105749, 37~pages,
 \href{https://arxiv.org/abs/2105.05660}{arXiv:2105.05660}.

\bibitem{BVM}
Bringmann K., Milas A., van Ittersum J.W.M., Quasimodularity  of $q$-series associated to root lattices, {i}n preparation.

\bibitem{CLM}
Calinescu C., Lepowsky J., Milas A., Vertex-algebraic structure of the
 principal subspaces of level one modules for the untwisted affine {L}ie
 algebras of types {$A$, $D$, $E$},
 \href{https://doi.org/10.1016/j.jalgebra.2009.09.029}{\textit{J.~Algebra}}
 \textbf{323} (2010), 167--192,
 \href{https://arxiv.org/abs/0908.4054}{arXiv:0908.4054}.

\bibitem{Carlsson}
Carlsson E., Vertex operators, {G}rassmannians, and {H}ilbert schemes,
 \href{https://doi.org/10.1007/s00220-010-1123-7}{\textit{Comm. Math. Phys.}}
 \textbf{300} (2010), 599--613,
 \href{https://arxiv.org/abs/0910.5528}{arXiv:0910.5528}.

\bibitem{DMN}
Dong C., Mason G., Nagatomo K., Quasi-modular forms and trace functions
 associated to free boson and lattice vertex operator algebras,
 \href{https://doi.org/10.1155/S1073792801000204}{\textit{Int. Math. Res.
 Not.}} \textbf{2001} (2001), 409--427,
 \href{https://arxiv.org/abs/math/0004132}{arXiv:math/0004132}.

\bibitem{DFR}
Dotsenko V., Feigin E., Reineke M., Koszul algebras and {D}onaldson--{T}homas
 invariants, \href{https://doi.org/10.1007/s11005-022-01604-4}{\textit{Lett.
 Math. Phys.}} \textbf{112} (2022), 106, 39~pages,
 \href{https://arxiv.org/abs/2111.07588}{arXiv:2111.07588}.

\bibitem{DM}
Dotsenko V., Mozgovoy S., D{T} invariants from vertex algebras,
 \href{https://doi.org/10.1017/S1474748024000288}{\textit{J.~Inst. Math.
 Jussieu}} \textbf{24} (2025), 291--339,
 \href{https://arxiv.org/abs/2108.10338}{arXiv:2108.10338}.

\bibitem{Efimov}
Efimov A.I., Cohomological {H}all algebra of a symmetric quiver,
 \href{https://doi.org/10.1112/S0010437X12000152}{\textit{Compos. Math.}}
 \textbf{148} (2012), 1133--1146,
 \href{https://arxiv.org/abs/1103.2736}{arXiv:1103.2736}.

\bibitem{Ekholm}
Ekholm T., Gruen A., Gukov S., Kucharski P., Park S., Sto\v{s}i\'{c} M.,
 Su{\l}kowski P., Branches, quivers, and ideals for knot complements,
 \href{https://doi.org/10.1016/j.geomphys.2022.104520}{\textit{J.~Geom.
 Phys.}} \textbf{177} (2022), 104520, 75~pages,
 \href{https://arxiv.org/abs/2110.13768}{arXiv:2110.13768}.

\bibitem{FS}
Feigin B., Stoyanovsky A.V., Quasi-particles models for the representations of
 {L}ie algebras and geometry of flag manifold,
 \href{https://arxiv.org/abs/hep-th/9308079}{arXiv:hep-th/9308079}.

\bibitem{Goncharov}
Goncharov A.B., Multiple polylogarithms, cyclotomy and modular complexes,
 \href{https://doi.org/10.4310/MRL.1998.v5.n4.a7}{\textit{Math. Res. Lett.}}
 \textbf{5} (1998), 497--516.

\bibitem{GM}
Goujard E., M\"{o}ller M., Counting {F}eynman-like graphs: {Q}uasimodularity
 and {S}iegel--{V}eech weight,
 \href{https://doi.org/10.4171/jems/924}{\textit{J.~Eur. Math. Soc. (JEMS)}}
 \textbf{22} (2020), 365--412,
 \href{https://arxiv.org/abs/1609.01658}{arXiv:1609.01658}.

\bibitem{CM}
Herzog J., Hibi T., Distributive lattices, bipartite graphs and {A}lexander
 duality,
 \href{https://doi.org/10.1007/s10801-005-4528-1}{\textit{J.~Algebraic
 Combin.}} \textbf{22} (2005), 289--302,
 \href{https://arxiv.org/abs/math/0307235}{arXiv:math/0307235}.

\bibitem{Hoffman}
Hoffman M.E., Algebraic aspects of multiple zeta values, in Zeta Functions,
 Topology and Quantum Physics, \textit{Dev. Math.}, Vol.~14,
 \href{https://doi.org/10.1007/0-387-24981-8_4}{Springer}, New York, 2005,
 51--73, \href{https://arxiv.org/abs/math/0309425}{arXiv:math/0309425}.

\bibitem{Zagier1}
Ihara K., Kaneko M., Zagier D., Derivation and double shuffle relations for
 multiple zeta values,
 \href{https://doi.org/10.1112/S0010437X0500182X}{\textit{Compos. Math.}}
 \textbf{142} (2006), 307--338.

\bibitem{JM1}
Jennings-Shaffer C., Milas A., On {$q$}-series identities for false theta
 series, \href{https://doi.org/10.1016/j.aim.2020.107411}{\textit{Adv. Math.}}
 \textbf{375} (2020), 107411, 22~pages,
 \href{https://arxiv.org/abs/2001.11368}{arXiv:2001.11368}.

\bibitem{JM2}
Jennings-Shaffer C., Milas A., Further {$q$}-series identities and conjectures
 relating false theta functions and characters, in Lie groups, Number Theory,
 and Vertex Algebras, \textit{Contemp. Math.}, Vol.~768,
 \href{https://doi.org/10.1090/conm/768/15467}{American Mathematical Society},
 Providence, RI, 2021, 253--269,
 \href{https://arxiv.org/abs/2005.13620}{arXiv:2005.13620}.

\bibitem{Kawasetsu}
Kawasetsu K., The free generalized vertex algebras and generalized principal
 subspaces,
 \href{https://doi.org/10.1016/j.jalgebra.2015.07.014}{\textit{J.~Algebra}}
 \textbf{444} (2015), 20--51,
 \href{https://arxiv.org/abs/1502.05276}{arXiv:1502.05276}.

\bibitem{Quivers}
Kucharski P., Reineke M., Sto\v{s}i\'{c} M., Su{\l}kowski P., Knots-quivers
 correspondence,
 \href{https://doi.org/10.4310/atmp.2019.v23.n7.a4}{\textit{Adv. Theor. Math.
 Phys.}} \textbf{23} (2019), 1849--1902,
 \href{https://arxiv.org/abs/1707.04017}{arXiv:1707.04017}.

\bibitem{Lang}
Lang S., Elliptic functions, 2nd ed., \textit{Grad. Texts in Math.}, Vol.~112,
 \href{https://doi.org/10.1007/978-1-4612-4752-4}{Springer}, New York, 1987.

\bibitem{Hao}
Li H., Some remarks on associated varieties of vertex operator superalgebras,
 \href{https://doi.org/10.1007/s40879-021-00477-6}{\textit{Eur.~J. Math.}}
 \textbf{7} (2021), 1689--1728,
 \href{https://arxiv.org/abs/2007.04522}{arXiv:2007.04522}.

\bibitem{HM}
Li H., Milas A., Jet schemes, quantum dilogarithm and {F}eigin--{S}toyanovsky's
 principal subspaces,
 \href{https://doi.org/10.1016/j.jalgebra.2023.10.033}{\textit{J.~Algebra}}
 \textbf{640} (2024), 21--58,
 \href{https://arxiv.org/abs/2010.02143}{arXiv:2010.02143}.

\bibitem{Libgober}
Libgober A., Elliptic genera, real algebraic varieties and quasi-{J}acobi
 forms, in Topology of Stratified Spaces, \textit{Math. Sci. Res. Inst.
 Publ.}, Vol.~58, Cambridge University Press, Cambridge, 2011, 95--120,
 \href{https://arxiv.org/abs/0904.1026}{arXiv:0904.1026}.

\bibitem{M2}
Milas A., Formal differential operators, vertex operator algebras and
 zeta-values,~{II},
 \href{https://doi.org/10.1016/S0022-4049(03)00036-7}{\textit{J.~Pure Appl.
 Algebra}} \textbf{183} (2003), 191--244,
 \href{https://arxiv.org/abs/math/0303154}{arXiv:math/0303154}.

\bibitem{M1}
Milas A., On certain automorphic forms associated to rational vertex operator
 algebras, in Moonshine: The First Quarter Century and Beyond, \textit{London
 Math. Soc. Lecture Note Ser.}, Vol.~372,
 \href{https://doi.org/10.1017/CBO9780511730054.016}{Cambridge University
 Press}, Cambridge, 2010, 330--357.

\bibitem{MP}
Milas A., Penn M., Lattice vertex algebras and combinatorial bases: general
 case and {$\mathcal W$}-algebras, \textit{New York J. Math.} \textbf{18}
 (2012), 621--650.

\bibitem{Ob-Serre}
Oberdieck G., A {S}erre derivative for even weight {J}acobi forms,
 \href{https://arxiv.org/abs/1209.5628}{arXiv:1209.5628}.

\bibitem{Ob}
Oberdieck G., Gromov--{W}itten invariants of the {H}ilbert schemes of points of
 a {K}3 surface, \href{https://doi.org/10.2140/gt.2018.22.323}{\textit{Geom.
 Topol.}} \textbf{22} (2018), 323--437,
 \href{https://arxiv.org/abs/1406.1139}{arXiv:1406.1139}.

\bibitem{Ohno}
Ohno Y., Okuda J., Zudilin W., Cyclic {$q$}-{MZSV} sum,
 \href{https://doi.org/10.1016/j.jnt.2011.08.001}{\textit{J.~Number Theory}}
 \textbf{132} (2012), 144--155.

\bibitem{Zagier2}
Ohno Y., Zagier D., Multiple zeta values of fixed weight, depth, and height,
 \href{https://doi.org/10.1016/S0019-3577(01)80037-9}{\textit{Indag.
 Math.~(N.S.)}} \textbf{12} (2001), 483--487.

\bibitem{Okounkov}
Okunkov A., Hilbert schemes and multiple {$q$}-zeta values,
 \href{https://doi.org/10.1007/s10688-014-0054-z}{\textit{Funct. Anal. Appl.}}
 \textbf{48} (2014), 138--144,
 \href{https://arxiv.org/abs/1404.3873}{arXiv:1404.3873}.

\bibitem{OEIS}
The on-line encyclopedia of integer sequences, \url{https://oeis.org/}.

\bibitem{Pan}
Pan Y., Peelaers W., Exact {S}chur index in closed form,
 \href{https://doi.org/10.1103/physrevd.106.045017}{\textit{Phys. Rev.~D}}
 \textbf{106} (2022), 045017, 34~pages,
 \href{https://arxiv.org/abs/2112.09705}{arXiv:2112.09705}.

\bibitem{Qin}
Qin Z., Yu F., On {O}kounkov's conjecture connecting {H}ilbert schemes of
 points and multiple {$q$}-zeta values,
 \href{https://doi.org/10.1093/imrn/rnw244}{\textit{Int. Math. Res. Not.}}
 \textbf{2018} (2018), 321--361,
 \href{https://arxiv.org/abs/1510.00837}{arXiv:1510.00837}.

\bibitem{Edge}
Renteln P., The {H}ilbert series of the face ring of a flag complex,
 \href{https://doi.org/10.1007/s003730200045}{\textit{Graphs Combin.}}
 \textbf{18} (2002), 605--619.

\bibitem{Ittersum}
van Ittersum J.W.M., Partitions and quasimodular forms: {V}ariations on the
 {B}loch--{O}kounkov theorem, Ph.D. thesis, {U}trecht {U}niversity, 2021,
 \url{https://staff.fnwi.uva.nl/j.w.m.vanittersum/files/theses/thesis_vanIttersum_withcover.pdf}.

\bibitem{Zagier}
Zagier D., Multiple zeta values, {P}reprint, 1995.

\bibitem{Zhu}
Zhu Y., Modular invariance of characters of vertex operator algebras,
 \href{https://doi.org/10.1090/S0894-0347-96-00182-8}{\textit{J.~Amer. Math.
 Soc.}} \textbf{9} (1996), 237--302.

\bibitem{Zudilin2}
Zudilin W., Algebraic relations for multiple zeta values,
 \href{https://doi.org/10.1070/RM2003v058n01ABEH000592}{\textit{Russian Math.
 Surveys}} \textbf{58} (2003), 1--29.

\bibitem{Zudilin1}
Zudilin W., Multiple $q$-zeta brackets,
 \href{https://doi.org/10.3390/math3010119}{\textit{Mathematics}} \textbf{3}
 (2015), 119--130, \href{https://arxiv.org/abs/1412.0163}{arXiv:1412.0163}.

\end{thebibliography}
\end{document}
